% =================================================================
% CONFIGURAÇÕES DO TEMPLATE - VARIANTE MONOGRAFIA
% Configuração completa para monografias e dissertações.
% Como usar:
%  - Preencha os placeholders abaixo (título, autor, etc.)
%  - Adicione/remova pacotes conforme a necessidade do seu trabalho
%  - Mantenha overrides aqui, sem editar config/sbc-template.sty
% Características desta variante:
%  - Suporte completo a capítulos (usar com documentclass report)
%  - Índice, lista de figuras, lista de tabelas
%  - Glossários e índice remissivo (opcionais)
%  - Todos os pacotes disponíveis
% =================================================================

% Pacotes essenciais
\usepackage{config/sbc-template}
\usepackage[T1]{fontenc}   % melhor hifenização/acentuação com pdflatex
\usepackage[utf8]{inputenc}
\usepackage{microtype}     % microtipografia (kerning/hifenização)
\usepackage[english,brazil]{babel} % idiomas: en + pt-BR
\usepackage{csquotes}      % citações tipográficas
\usepackage{graphicx,url}
\usepackage{hyperref}
\graphicspath{{assets/}}   % pasta padrão para imagens

% Pacotes adicionais (todos incluídos para monografias)
\usepackage{listings}           % listagens de código
\usepackage{xcolor}             % cores (necessário para listings)
\usepackage{amsmath,amssymb}    % matemática
\usepackage{booktabs}           % tabelas profissionais (\toprule, etc.)
\usepackage{tabularx}           % tabelas largas com colunas X
\usepackage{array}              % colunas avançadas em tabelas
\usepackage{enumerate}          % listas enumeradas flexíveis

% Glossário e índice (opcionais - descomente se necessário):
% \usepackage[acronym]{glossaries} % requer makeglossaries; ver README
% \makeglossaries
% \usepackage{imakeidx}           % índice remissivo
% \makeindex

% BibLaTeX (opção B2 - descomente se preferir ao invés de BibTeX):
% \usepackage[backend=biber,style=authoryear,doi=true,isbn=false,url=true]{biblatex}
% \addbibresource{bibliography/sbc-template.bib}

% Minted (opcional, para código com realce de sintaxe via Pygments):
% Requer: \usepackage{minted} e compilar com TEXFLAGS="-shell-escape".
% \usepackage{minted}

% Estilo de código (exemplo; ajuste ou remova)
\definecolor{codegray}{RGB}{248,248,248}
\lstdefinestyle{code}{
    backgroundcolor=\color{codegray},
    basicstyle=\footnotesize\ttfamily,
    numbers=left,
    numbersep=5pt,
    frame=single,
}

% Configurações de formatação para monografias
% Espaçamento entre capítulos no sumário
\usepackage{tocloft}
\setlength{\cftbeforechapskip}{0.5em}

% Comandos de placeholders (consumidos em \title, \author e \address)
\newcommand{\DocumentTitle}{My Monograph}
\newcommand{\AuthorName}{Jane Doe}
\newcommand{\AuthorInstitution}{Research Institute}
\newcommand{\AuthorLocation}{New York -- NY -- USA}
\newcommand{\AuthorEmail}{jane@example.com}
\newcommand{\AdvisorName}{[Nome do Orientador]}
\newcommand{\AdvisorTitle}{[Prof. Dr./Dra.]}

% Hyperlinks (preenche metadados do PDF)
\hypersetup{
    colorlinks=true,
    linkcolor=black,
    urlcolor=blue,
    citecolor=blue,
    pdftitle={\DocumentTitle},
    pdfauthor={\AuthorName},
    pdfsubject={[Área de Estudo]},
    pdfkeywords={[palavras-chave, separadas, por vírgula]},
    bookmarks=true,
    bookmarksopen=true,
    bookmarksnumbered=true,
}

% Metadados padrão (para projetos que usam este arquivo como entrada principal)
\title{\DocumentTitle}
\author{\AuthorName\inst{1}}
\address{\AuthorInstitution\\
  \AuthorLocation\\
  \email{\AuthorEmail}
}

% Configurações gerais
\sloppy
\setlength{\parskip}{6pt}
